\documentclass[11pt]{article}
\usepackage[margin=1in]{geometry}
\usepackage{amsmath,amssymb,booktabs,hyperref}
\newcommand{\F}{\mathbb F}
\newcommand{\E}{\mathbb E}
\title{Counting pooled determinant zeros from planted relations and field moments}
\author{Project research note}
\date{20 November 2026}
\begin{document}
\maketitle

Nine pooled field moments give an exact zero-count certificate on a
testable subset of three-slot inputs with two abscissa triples. A learner
trained on planted relations resolves some of the remaining ambiguities.
At the growing arity of the proposed reduction, the requisite fresh error
and end-to-end exponent saving remain open.

\section{A zero-count certificate using two hidden centers}
Use the independent dyadic hash windows and the unweighted determinant
$\Delta$ of \texttt{target\_incidence.tex}. Write $F_i$ for the retained
affine points, $D=\prod|F_i|$, $H$ for the number of ordered tuples with
sum the prescribed target $T$, and $Z_j=\sum_{\mathbf P}\Delta_{\mathbf P}^j$.
Regularity means disjoint abscissas between slots and avoidance of $x(T)$.
It guarantees that the multiplicity of zero among the determinants is $H$.
Here $m=3$, $p>7$, and the three regular slots have respectively one, one,
and two distinct signed pairs in any order. This shape can be tested while
constructing the slots. The decoder subsequently uses only $Z_1,\ldots,Z_k$.
Fixing an abscissa in each slot leaves eight choices of signs. The determinant
in the basis $1,x,x^2,y$ is affine in the three signs because the determinant
expansion uses the $y$ row once. Thus there are two equally weighted blocks
of the form $c_\sigma+\sum_{i=1}^3 a_{\sigma i}\epsilon_i$, $\epsilon_i=\pm1$.

Set $c=Z_1/16$ and denote the block centers by $c\pm h$. Let $M_j$ be the
mean of $(\Delta-c)^j$, computed by the binomial formula from the raw
moments, with $M_0=1$. Write $t=h^2$, $A_0=1$, $B_0=0$. For $q=1,2,3$ put
\[
 A_q=M_{2q}-\sum_{r<q}{2q\choose 2r}t^{q-r}A_r,\qquad
 B_q=\frac{M_{2q+1}-\sum_{r<q}{2q+1\choose 2r}t^{q-r}B_r}{2q+1}.
\]
Expanding around each center and using the vanishing of odd centered
moments gives block means $A_q\pm B_q/h$ for the local centered $2q$th
powers whenever $h\ne0$. In particular, seven moments determine these
means if the distinct centers are supplied. For hidden centers, do
arithmetic over $\F_p[h]$ and scale the determinant by $h$. Then the
local means for the scaled centered values are polynomial:
\[
 u_{0,\sigma}=1,\qquad u_{2q,\sigma}=h^{2q}A_q+\sigma h^{2q-1}B_q .
\]
Within a block omit the index $\sigma$. The first six local moments give
\[
 s_1=u_2,\quad s_2=(3s_1^2-u_4)/2,\quad
 s_3=(u_6-15u_4s_1+30s_1^3)/16,\quad
 e_2=(s_1^2-s_2)/2,\quad e_3=(s_1^3-3s_1s_2+2s_3)/6 .
\]
Writing $a_i$ for the scaled sign coefficients, $s_j=\sum a_i^{2j}$
and $e_j$ are elementary symmetric functions of the $a_i^2$.
The block octic from \texttt{full\_moment\_decoding.tex} is
\[
 K_\sigma(z)=((z^2-s_1)^2-4e_2)^2-64z^2e_3 .
\]
Its roots with multiplicity are the eight centered scaled sign values.
Use its recurrence to predict $u_{8,\sigma},u_{10,\sigma},\ldots$ and take
the gcd of the nonzero residuals for $8\le j\le k$,
\[
 R_j(h)=h^jM_j-\tfrac12\sum_{\sigma=\pm1}
  \sum_{r=0}^{\lfloor j/2\rfloor}{j\choose 2r}
  (\sigma h^2)^{j-2r}u_{2r,\sigma}.
\]
Every genuine nonzero separation is a root of each residual. If all
available odd $M_j$ vanish, an equal-center case remains possible and the
implemented certificate abstains. It also abstains if every residual
vanishes identically. Otherwise intersect the gcd with $h^{p-1}-1$ using
modular powering. The resulting polynomial $G$ has simple roots exactly
at the surviving nonzero field candidates.

Only the first five coefficients at $z=0$ of
\[
 K_+(z-hc-h^2)\,K_-(z-hc+h^2)
\]
are needed. The two smaller slots have four ordered pairs and each pair
has at most one completing point, so $H\le4$. For $r=0,\ldots,4$ intersect
the remaining $G$ with its gcd with coefficient $[z^r]$. A degree drop
records a candidate of multiplicity $r$, and the gcd retains exactly the
candidates that still vanish to higher order. This gives a superset of
the possible counts containing the true count. Return the count if all
agree, or the existence bit if they agree on positivity. Spurious candidates
can lead to abstention. The constant coefficient alone tests vanishing
at a given separation. We never recover the field roots of a block, the
individual centers, or the degree-16 polynomial of the observed spectrum.
For a fixed prefix, e.g. $k=9$, this readout uses bounded-degree gcds
and $O(\log p)$ field operations including modular powering. Indeed
$\deg_h u_{2r,\sigma}\le4r$ by the stated construction and octic
recurrence, hence $\deg_hR_j\le2j$. Computing the moment prefix from the
seeds is a separate charge as in \texttt{target\_incidence.tex}.

\section{Planted learning and exact finite checks}
Write $P$ for fresh seeds and $Q$ for seeds conditioned to retain a
uniformly selected affine tuple of sum $T$. Sample all but one point
uniformly, force the remaining point by the sum, retry if it is the
identity, and sample each slot hash conditionally on its own point.
Then $dQ/dP=H/\mu$, where $\mu=\E_P H$. On an ambiguous bin $v$,
the score $2\mu Q(v)>P(v)$ predicts presence. It compares the
conditional mean hit count with $1/2$. Its difference from optimal
existence scoring is controlled in fresh risk by the mass of extra
witnesses, $O(\lambda^2)$ at the asymptotic tuning. On the small
fixtures we report the actual error.

For $y^2=x^3+4x+1$ over $\F_{13}$ (group order 19), all odd multipliers
modulo 16, all offsets, and widths $(2,2,2)$ give 168,480 eligible seed
triples counting multiplicities. Of these 133,008 contain a relation.
The following numbers are exact on that conditional population. The
independent reference checks the zero count against the elliptic sum.
\begin{center}
\begin{tabular}{rrrr}
\toprule Moments & Count certified & Existence certified & Misses if abstentions mean absent\\
\midrule
8 & 123312 & 123792 & 36096\\
9 & 145776 & 146256 & 13632\\
10 & 151008 & 151008 & 10272\\
12 & 151008 & 151008 & 10272\\
\bottomrule
\end{tabular}
\end{center}
Next, keep all nine-moment certificates and group unresolved inputs by
the candidate-count set, including a bin for abstention. Use independent
draws of 60,000 plants and 60,000 unlabeled fresh seeds at the fixed
target to fit the score. No ground-truth fresh label is used to train the
decisions. Two bins were called positive. Exact evaluation on the whole
eligible fresh population gave 7,680 errors, of which 7,200 were false
positives and 480 were false negatives. This improves conditional existence
error from 8.09\% to 4.56\% relative to rejecting every ambiguous input.
Positive decisions in ambiguous bins do not certify multiplicity.

We also trained a rule for the \emph{integer} count. On ambiguous training
plants only, count the completions of the four small-slot pairs by group
addition and membership in the remaining mask. For $j=1,\ldots,4$ the
population identity is
\[
 P(H=j,F=v)=\frac{\mu}{j} Q(H=j,F=v).
\]
Subtract their sum from the unlabeled bin mass $P(F=v)$ to estimate the
mass at zero. Predict the most frequent admissible count from these
estimated masses, retaining all certified counts. The same train draws
give 10,320 incorrect integer counts, compared with 14,112 using the
smallest admissible count (zero on abstention). Thus 158,160 of 168,480
counts are correct, including 145,776 certified ones. Inference uses
only the pooled moments and shape. The cheap count annotation of a
training plant uses the four-pair bound specific to this case.

Separately, on $y^2=x^3+2x+2$ over $\F_{17}$, $T=(0,6)$, the first nine
moments certified 375 exact counts and 376 bits on all 420 distinct
regular abscissa arrangements up to swap symmetries. These arrangements
have equal weight in this particular check.

\section{Larger pools and independent fresh seeds}
At fixed $T$, we fitted two families of field equations to genuine
positive examples. The first uses every monomial of weighted degree at
most eight in the second, fourth and sixth centered moments after
division of every determinant by its nonzero pooled mean. There are
11 such monomials. The second uses all 55 monomials of ordinary degree
at most two in the nine \emph{raw field values} $Z_1,\ldots,Z_9$,
including odd moments and inputs with vanishing mean. Each family fits
all equations in the planted nullspace. Sampling 12,000 independent
plants per curve restored the known one-orbit identity in the normalized
family over $\F_{101}$. On regular masks with at least three orbits the
planted feature matrices had full ranks 11 and 55 on both curves. Thus
every nonzero equation in either family fails on at least one genuine
plant in its tested stratum, an exact finite statement. Independent
samples conditioned on retaining independent affine marks in all slots
provided a control for the bias toward larger slots.

The same trials also bound \emph{approximate} vanishing on the planted
population. Partition the eligible plant samples into disjoint blocks of
$d$ rows, where $d$ is the number of features. If a hyperplane contains
mass $a$, each block is singular with probability at least $a^d$. This
holds for the hyperplane of largest mass and thus covers coefficients
chosen after seeing the samples. Conditional on the number of eligible
draws, the accepted rows have the independent conditional law. Invert
the one-sided binomial test at level $0.005$ on the block singularity
probability, then take the $d$th root. Under independent draws the
result is a 99.5\% upper confidence bound for the maximum fraction of
eligible plants captured by one nonzero equality in that family.
\begin{center}
\begin{tabular}{rrrrr}
\toprule Field & Features & Eligible plants & Singular blocks & Recall upper bound\\
\midrule
101 & 11 & 7254 & 4/659 & 0.6974\\
1031 & 11 & 10710 & 0/973 & 0.6224\\
101 & 55 & 7328 & 1/133 & 0.9485\\
1031 & 55 & 10716 & 0/194 & 0.9364\\
\bottomrule
\end{tabular}
\end{center}
The confidence levels apply separately to the four rows. The rates
describe specific finite input laws and the indicated families of
single equalities. They give no lower bound on unrestricted decoding.

We then tested 16,000 independently drawn fresh seeds per curve with
widths $(4,4,4)$ at $p=101$ and $(8,16,16)$ at $p=1031$. Both use
$y^2=x^3+ax+1$, respectively $a=17,9$, with group orders 103 and 1019,
and the first affine generator in the deterministic fixture as target.
We used the planted equations available for each stratum and the
nine-moment certificate on the two-orbit stratum. The rule returns
absent if it has no applicable test, including on irregular inputs.
Unconditionally this gave the following integer error counts.
\begin{center}
\begin{tabular}{rrrrr}
\toprule Field & Fresh positives & Two-orbit exact counts & Large-pool positives missed & All errors\\
\midrule
101 & 3669 & 1715/1732 & 2299/2299 & 3363\\
1031 & 3274 & 44/44 & 2929/2929 & 3272\\
\bottomrule
\end{tabular}
\end{center}
There were no false positives in these two trials. The calibrated
$\F_{13}$ ambiguity rule was trained for its own field and target.
Truth and the input moments were obtained by enumeration on small
fixtures. None of the claimed algorithmic costs use the enumerator.

\section{Why planting needs a genuinely joint observation}
There is a fixed-target obstruction for any rule that leaves even one
slot unread. It applies to arbitrary arithmetic on the observed seeds.
Let $R$ generate $m$ independent uniformly marked affine points, then
condition each hash on retaining its mark, as in the matched experiment.
Writing $P$ for fresh inputs gives $dR/dP=D/\E_P D$. Given the hashes
under $R$, the marks are independent uniform members of their slots.
Conditioning the marks' sum to $T$ gives exactly $Q$.

For a nonempty proper subset $I$ of $s$ slots, let $t=m-s$,
$D_I=\prod_{i\in I}|F_i|$, and let $H_I$ count partial tuples of sum $T$.
Under $R$, $D_I>0$ almost surely. If $a_j$ is the probability that $j$
independent uniform affine points sum to the nonidentity $T$, then
$a_j=(1-(-1)^j/(N-1)^j)/N$. More generally the probability of sum $u$
is $[1+(-1)^j(N1_{u=O}-1)/(N-1)^j]/N$, by expanding the $j$th
convolution power of the uniform nonidentity measure. Conditioning
on the observed hashes and averaging over their marks gives
\[
 \frac{dQ_I}{dR_I}=
 \frac{1+(-1)^t(NH_I/D_I-1)/(N-1)^t}{1-(-1)^m/(N-1)^m}.
\]
Since $X=H_I/D_I$ lies in $[0,1]$ and $\E_R X=a_s$,
$\E_R|X-a_s|\le 2a_s(1-a_s)$ by convexity. Consequently
\[
 \|Q_I-R_I\|_{\rm TV}\le
 \frac{Na_s(1-a_s)}{(N-1)^t(1-(-1)^m/(N-1)^m)}
 =O((N-1)^{-t}).
\]
A randomized rule may adaptively select its next slot from its previous
observations. If it reads at most $m-1$ slots on every path, partition
an acceptance event by the first omitted index in numerical order.
The part for index $i$ is measurable on all slots except $i$: simulate
until the decision or the first query to $i$, which excludes that part.
Summing the bounds and averaging over the rule's independent random tape
gives $|Q(A)-R(A)|=O(m/N)$. The product-variance bound in the final
section also gives $\|R-P\|_{\rm TV}\le
\tfrac12\sqrt{\operatorname{Var}_P D}/\E_P D=o(\lambda)$.
Thus $Q(A)\le P(A)+o(\lambda)$. Its fresh existence error obeys
\[
 P(Y=1)+P(A)-2P(Y=1,A)
 \ \ge\ P(Y=1)+(1-2\mu)P(A)-o(\lambda^2)
 \ \ge\ \lambda-O(\lambda^2),
 \qquad Y=1_{\{H>0\}}.
\]
The first inequality uses $P(Y=1,A)\le\mu Q(A)$. The same lower bound
holds for the exact count by turning any count output into an existence
bit. A single pooled $Z_j$ already depends on all
slots and falls outside the lower bound. Providing any additional
information from omitted slots would also change the model.
On the $\F_{13}$ fixture, direct summation over actual hash
multiplicities verified the marginal likelihood identity on all 756
nonempty one- and two-slot mask patterns in the tested positions.

\section{What a planted dictionary would have to learn}
For any feature $F$ of the full moments, including the entire moment
vector, let $\tau_s=\min_{|K|\le s}Q(F\notin K)$. A dictionary that
predicts presence on at most $s$ feature values can achieve fresh error
$O(\lambda^2)$ only if $\tau_s=O(\lambda)$, since
\[
 \mu Q(F\notin K)\le P(H>0,F\notin K)+\E_P(H-1)_+ .
\]
Conversely, suppose this concentration holds and the optimal classifier
given $F$ has fresh error $O(\lambda^2)$. Draw $n=\lceil s/\lambda\rceil$
plants and retain their feature keys. The expected unseen planted mass is
at most $\tau_s+s/(en)$ since $q(1-q)^n\le1/(en)$. Restricting the optimal
positive decisions to the discovered keys increases expected fresh error
by at most $\mu$ times this missing mass. Use independent planted and
fresh draws to fit the bin decisions as above. Conditional on discovery,
the class consists of at most $2^n$ subsets. Hoeffding and a union bound
give uniform additive score error $O(\lambda^2)$ with failure probability
$O(\lambda^2)$ using $O((n+\log(1/\lambda))/\lambda^4)$ draws of each
kind. Per-bin counters implement the empirical minimum in time linear
in the sample count up to feature evaluation and lookup costs.
Thus a positive dictionary of size $2^{O(b/m)}$ would suffice on this
exponential scale if the remaining conditions held, even with an
enormous negative alphabet. Concentration alone says nothing about
collisions between negative and positive seeds.

\section{End-to-end accounting}
Set $b=\lceil\log_2p\rceil$, $L=\log_2b$, $g=\log_2 L$, with
$1/(2L)<\lambda\le1/L$ and $m\asymp L^{1/4}\sqrt g$. The inherited
estimates give $\mu\sim\lambda$ and $\E_P H(H-1)=O(\lambda^2)$ at
each prescribed target. Hence a decoder meeting the requested error
must accept a $1-O(\lambda)$ fraction of plants while assigning
$O(\lambda^2)$ fresh mass to false positives. It has a budget of
$2^{O(b/m)}$ bit operations inclusive of sampling, seed-dependent
setup, moment evaluation, and decoding. A cheap moment prefix alone
does not satisfy that requirement.

Typically the number of signed abscissa blocks is $D/2^m$. Indeed the
probability of retaining any affine point of order two is at most
$3\sum_i\rho_i=o(\lambda^2)$, where $\rho_i$ is the slot retention rate.
Slot independence and $\operatorname{Var}|F_i|\le2\E|F_i|$ yield
$\operatorname{Var}D/(\E D)^2\le\prod_i(1+2/\E|F_i|)-1$.
This is $o(\lambda^2)$ at balanced widths, and $\E D=(1+o(1))N\lambda$
for $N=|E(\F_p)|$. With probability $1-o(\lambda^2)$ the number of
blocks thus lies within constant factors of $N\lambda/2^m=2^{b-o(b)}$.
Explicitly visiting them in pairs already exceeds the budget. Also,
the higher-arity determinant can contain products of sign variables.
A fixed-target decoder operating directly on the pooled information
for these inputs remains open. Under packet assumption A.1 and the
inherited contraction and reduction estimates the best complete bound
in the project remains
\[
 \log_2 \E T_{\rm solve}
 \le (26.2933775156+o(1))\,b/(\log_2 b)^{1/4}.
\]

The implementation is \texttt{pooled\_zero.py}. From the project root run
\texttt{python -B research/joint\_parity/check\_pooled\_zero.py} for
the exact checks and planted calibration, and
\texttt{python -B research/joint\_parity/probe\_pooled\_zero.py}
for the independent trials. The scripts fix their pseudorandom seeds.
\end{document}
